% ==============================================================================
% ================================ AUFGABE 3 ===================================
% ==============================================================================

%\chapter{Aufgabe 3: KV-Diagramm}
%\label{chap:aufgabe3}
%\thispagestyle{fancy}

\section{Aufgabe 3: KV-Diagramm}
In diesem Kapitel wird die gegebene Wertetabelle für die vier Ausgänge \(Q_0, Q_1, Q_2, Q_3\) in Abhängigkeit von den vier Eingangssignalen \(A, B, C, D\) ausgewertet. 
Es werden die KV‑Diagramme aller Ausgänge gezeichnet, die zugehörigen minimierten logischen Gleichungen ermittelt und abschließend die Gesamtschaltung skizziert, so wie in dem Studienheft \cite{2} beschrieben.
\subsection{Wertetabelle}
Die Tabelle enthält die 16 möglichen Eingangskombinationen (\textit{DCBA}) und die zugehörigen Ausgänge \(Q_3 Q_2 Q_1 Q_0\) gemäß der Aufgabenstellung.
\begin{table}[h]
\centering
\begin{tabular}{c|c|c|c|c|c|c|c}
\toprule
\(D\) & \(C\) & \(B\) & \(A\) & \(Q_3\) & \(Q_2\) & \(Q_1\) & \(Q_0\) \\
\midrule
0 & 0 & 0 & 0 & 0 & 0 & 0 & 0 \\
0 & 0 & 0 & 1 & 0 & 0 & 0 & 1 \\
0 & 0 & 1 & 0 & 0 & 0 & 1 & 1 \\
0 & 0 & 1 & 1 & 0 & 0 & 1 & 0 \\
0 & 1 & 0 & 0 & 0 & 1 & 1 & 0 \\
0 & 1 & 0 & 1 & 0 & 1 & 1 & 1 \\
0 & 1 & 1 & 0 & 0 & 1 & 0 & 1 \\
0 & 1 & 1 & 1 & 0 & 1 & 0 & 0 \\
1 & 0 & 0 & 0 & 1 & 1 & 0 & 0 \\
1 & 0 & 0 & 1 & 1 & 1 & 0 & 1 \\
1 & 0 & 1 & 0 & 1 & 1 & 1 & 1 \\
1 & 0 & 1 & 1 & 1 & 1 & 1 & 0 \\
1 & 1 & 0 & 0 & 1 & 0 & 1 & 0 \\
1 & 1 & 0 & 1 & 1 & 0 & 1 & 1 \\
1 & 1 & 1 & 0 & 1 & 0 & 0 & 1 \\
1 & 1 & 1 & 1 & 1 & 0 & 0 & 0 \\
\bottomrule
\end{tabular}
\caption{Wertetabelle für \(Q_0 \dots Q_3\) (\textit{Quelle: Eigene Darstellung})}
\label{tab:wahrheit}
\end{table}
\subsection{KV-Diagramme und minimierte Gleichungen}
Für jeden Ausgang wird ein KV‑Diagramm mit der Anordnung \(AB\) (\textit{horizontal}) und \(CD\) (\textit{vertikal}) im Gray‑Code \((\textit{00}, \textit{01}, \textit{11}, \textit{10})\) erstellt. 
Die Einsen werden zu möglichst großen Blöcken zusammengefasst.
\subsubsection{\texorpdfstring{Ausgang \(Q_0\)}{Ausgang Q0}}
\begin{table}[h]
\centering
\begin{tabular}{c|cccc}
\multicolumn{1}{c}{} & \multicolumn{4}{c}{\(AB\)} \\
\cline{2-5}
\(CD\) & 00 & 01 & 11 & 10 \\
\hline
00 & 0 & 1 & 0 & 1 \\
01 & 0 & 1 & 0 & 1 \\
11 & 0 & 1 & 0 & 1 \\
10 & 0 & 1 & 0 & 1 \\
\hline
\end{tabular}
\caption{KV-Diagramm für \(Q_0\) (\textit{Quelle: Eigene Darstellung})}
\end{table}
Die minimierte Gleichung lautet: \hspace{4cm} \[\boxed{Q_0 = \overline{A} \cdot B \;\cup\; A \cdot \overline{B}}\]
\subsubsection{\texorpdfstring{Ausgang \(Q_1\)}{Ausgang Q1}}
\begin{table}[h]
\centering
\begin{tabular}{c|cccc}
\multicolumn{1}{c}{} & \multicolumn{4}{c}{\(AB\)} \\
\cline{2-5}
\(CD\) & 00 & 01 & 11 & 10 \\
\hline
00 & 0 & 0 & 1 & 1 \\
01 & 1 & 1 & 0 & 0 \\
11 & 1 & 1 & 0 & 0 \\
10 & 0 & 0 & 1 & 1 \\
\hline
\end{tabular}
\caption{KV-Diagramm für \(Q_1\) (\textit{Quelle: Eigene Darstellung})}
\end{table}
Die minimierte Gleichung lautet: \hspace{4cm} \[\boxed{Q_1 = A \cdot \overline{C} \;\cup\; \overline{A} \cdot C}\]
\subsubsection{\texorpdfstring{Ausgang \(Q_2\)}{Ausgang Q2}}
\begin{table}[h]
\centering
\begin{tabular}{c|cccc}
\multicolumn{1}{c}{} & \multicolumn{4}{c}{\(AB\)} \\
\cline{2-5}
\(CD\) & 00 & 01 & 11 & 10 \\
\hline
00 & 0 & 0 & 0 & 0 \\
01 & 1 & 1 & 1 & 1 \\
11 & 0 & 0 & 0 & 0 \\
10 & 1 & 1 & 1 & 1 \\
\hline
\end{tabular}
\caption{KV-Diagramm für \(Q_2\) (\textit{Quelle: Eigene Darstellung})}
\end{table}
Die minimierte Gleichung lautet: \hspace{4cm} \[\boxed{Q_2 = \overline{C} \cdot D \;\cup\; C \cdot \overline{D}}\]
\subsubsection{\texorpdfstring{Ausgang \(Q_3\)}{Ausgang Q3}}
\begin{table}[h]
\centering
\begin{tabular}{c|cccc}
\multicolumn{1}{c}{} & \multicolumn{4}{c}{\(AB\)} \\
\cline{2-5}
\(CD\) & 00 & 01 & 11 & 10 \\
\hline
00 & 0 & 0 & 0 & 0 \\
01 & 1 & 1 & 1 & 1 \\
11 & 1 & 1 & 1 & 1 \\
10 & 0 & 0 & 0 & 0 \\
\hline
\end{tabular}
\caption{KV-Diagramm für \(Q_3\) (\textit{Quelle: Eigene Darstellung})}
\end{table}
Die minimierte Gleichung lautet: \hspace{4cm} \[\boxed{Q_3 = D}\]
\subsection{Gesamtschaltung}
Die folgende Abbildung zeigt die Gesamtschaltung aller Ein- und Ausgänge unter Verwendung der Logik-Bausteine UND, ODER und NEGATION (\textit{Inverter}).
\begin{figure}[h]
\centering
\includegraphics[width=0.44\textwidth]{asset/image/Aufgabe3.png}
\caption{Gesamtschaltung aller Ein- und Ausgänge mit UND-, ODER- und NICHT-Gattern (\textit{Quelle: Eigene Darstellung})}
\label{fig:aufgabe3_schaltung}
\end{figure}
Die Schaltung realisiert die minimierten Gleichungen: \hspace{4cm} \[\boxed{
\begin{aligned}
    Q_0 &= \overline{A} \cdot B \;\cup\; A \cdot \overline{B} \\
    Q_1 &= A \cdot \overline{C} \;\cup\; \overline{A} \cdot C \\
    Q_2 &= \overline{C} \cdot D \;\cup\; C \cdot \overline{D} \\
    Q_3 &= D
\end{aligned}
}\]

% ==============================================================================
% ================================ AUFGABE 3 ===================================
% ==============================================================================